\section{LLM 的贡献与边界：虚拟空间不是全部世界}

前面已经建立世界模型的定义，接下来要把它和 LLM 区分开。这个区分很重要，因为访谈并不是反 LLM，而是在问 LLM 的成功到底覆盖了世界的哪一部分。

谢赛宁并不否认 LLM 的革命性。他说自己要感谢 LLM，因为没有 LLM，多模态智能也不会扩展到今天的规模。但他反对把 LLM 叙事扩展成“语言模型自然通往 human-level intelligence”。在他的框架里，LLM 更擅长 digital / virtual space：文本、代码、知识总结、教育、法律、搜索和 agentic coding 等。它可以成为智能系统的重要元素，却不是世界模型的根基。

\begin{figure}[H]
\centering
\includegraphics[width=0.96\textwidth]{llm-vs-world-model.png}
\caption{LLM 与世界模型：两种智能市场的对比。}
\end{figure}

\begin{knowledgebox}{读图：语言模型和世界模型互补，但不能互相替代}
左侧 LLM 的强项是 token 化后的知识、文本推理和数字空间操作；右侧 world model 的目标是连续信号、物理预测、行动规划和安全控制。二者可以互补：语言提供接口、知识和沟通；世界模型提供物理接地、动态预测和行动能力。把一边强行归约成另一边，会误判技术瓶颈。
\end{knowledgebox}

\subsection{为什么语言模型像强监督学习}

访谈里有一个很有启发的观点：语言模型常被说成 self-supervised learning，但从另一种角度看，语言本身已经是人类文明加工后的强监督信号。几千年文明、书籍、网页、代码、论文和互联网，把大量世界知识压缩成 tokenized 的文本。训练 LLM 像是下载这些已加工知识，而不是直接学习物理世界。

这解释了为什么 LLM scaling law 能较早出现：它有大量已经被人类整理、标注、沟通化的材料。相反，vision 和 robotics 面对的是 raw sensory data、动作、空间、物理约束和反馈，天然缺少同等规模、同等质量、同等压缩程度的“互联网式标签”。

\begin{warningbox}{“免费数据”不等于“无监督数据”}
互联网上的文本可以免费抓取，并不意味着它没有人类标签。语言是人类为了沟通而长期压缩出来的结构；模型吃下这些文本时，也吃下了人类已经完成的大量抽象、选择和标注。
\end{warningbox}

\subsection{杯子摔碎：语言描述遗漏了什么}

谢赛宁举了一个简单例子：我们说“杯子掉在地上碎了”，语言只保留了沟通所需的信息。它没有描述杯子的接触、受力、破裂路径、材料性质、碎片飞散和声音。对人类沟通来说，这些细节大多没必要；对一个要在物理世界行动的系统来说，这些细节中的一部分可能非常关键。

所以，语言不是世界本身，而是为了交流而做的压缩。世界模型需要学习的，恰恰是语言压缩之外的那些规律：动力学、空间结构、接触、时序、因果和可行动性。

\subsection{本章小结}

LLM 的成功来自人类知识的巨大压缩和互联网规模。它会继续是智能系统的重要组成部分，但面向物理世界的智能不能只靠语言。世界模型要补的是语言天然省略掉的那部分世界。

